% ===================================================================
%  LENTELĖ Nr. 15 — Turgus. --- Maisto prekės.
%  Traduction 2026 d'apres texto/io, controlee sur texto/fr, texto/ru
%  et texto/cs. Consigne : voir texto/lt/10-lentele-01.tex.
%
%  LE MARCHE EST, POUR LE GALICIEN, L'UN DES ENDROITS OU LA DISTANCE
%  A LA LANGUE VOISINE EST MAXIMALE, et la colonne tcheque y avait
%  aligne cinq charcuteries sans un emprunt. LE LITUANIEN FAIT MIEUX
%  QUE CONFIRMER : il permet de DATER.
%
%  L'IRLANDAIS, AU TABLEAU 14 DE LA ONZIEME COLONNE, AVAIT DONNE LA
%  REGLE — la date de l'objet decide la forme du mot. L'etal des
%  legumes est le meilleur endroit du livret pour la verifier, parce
%  que les plantes ont des dates d'arrivee connues :
%
%    du fonds ou si vieux qu'on ne le sent plus : « ropė » le navet,
%      « kopūstas » le chou, « svogūnas » l'oignon, « česnakas »
%      l'ail, « žirnis » le pois, « pupa » la feve, « morka » la
%      carotte, « grybas » le champignon, « obuolys » la pomme,
%      « kriaušė » la poire, « slyva » la prune, « riešutas » la
%      noix ;
%    venus d'Amerique, et empruntes tous les trois : « bulvė » la
%      pomme de terre, « pomidoras » la tomate, « kukurūzai » le
%      mais ;
%    venus de plus loin encore, et empruntes de meme : « arbata » le
%      the, « kava » le cafe, « ryžiai » le riz, « cukrus » le
%      sucre, « bananai », « ananasai », « šokoladas ».
%
%  ET L'ETAL DU POISSONNIER CONFIRME LE TABLEAU 10. La Baltique a ses
%  poissons et le lituanien les nomme : « menkė » la morue,
%  « silkė » le hareng, « ungurys » l'anguille, « lynas » la tanche,
%  « karpis » la carpe, « upėtakis » la truite — mot a mot le
%  coureur-de-riviere —, « vėžys » l'ecrevisse. Mais la langouste,
%  le homard, l'huitre et la moule de l'ocean restent
%  « langustas », « omaras », « austrė », « midija ». UNE LANGUE
%  NOMME LES POISSONS DE SA MER ET EMPRUNTE CEUX DES AUTRES.
%
%  Le releve des viandes, lui, ne coute pas un emprunt, comme en
%  tcheque : « kumpis » le jambon, « dešra » la saucisse,
%  « kraujinė » le boudin, « lašiniai » le lard, « aviena »,
%  « jautiena », « veršiena » — les trois viandes formees sur le nom
%  de la bete avec le meme suffixe.
% ===================================================================

%%K t15-apar-1 apar
\VUsaut{4.0mm}
\VUtitre{15.0pt}{60}{LENTELĖ Nr. 15}
\VUsaut{3.0mm}

%%K t15-tit-1 sub
\VUpk{11.4pt}{40}{Turgus. --- Maisto prekės.}
\VUsaut{3.0mm}

%%K t15-01-1 p
1. --- Dauguma prekybininkų, kurie tiekia mums mūsų maistui reikalingas prekes,
renkasi \VUgras{dengto turgaus} \VUgras{halėje} \textsuperscript{(1)} arba gretimose
gatvėse.

%%K t15-02-1 p
2. --- \VUgras{Kiaulienos mėsininkas} \textsuperscript{(2)}, kuris parduoda
\VUgras{kumpį} \textsuperscript{(3)}, \VUgras{kraujinę dešrą} \textsuperscript{(4)},
\VUgras{dešrelę} \textsuperscript{(5)}, \VUgras{vidurių dešrelę} \textsuperscript{(6)} ir
\VUgras{lašinius} \textsuperscript{(7)}, stovi prie \VUgras{sviesto ir sūrio
pardavėjos} \textsuperscript{(8)}, kurios prekystalis aprūpintas
\VUgras{olandiškais sūriais} \textsuperscript{(9)} su raudona pluta,
\VUgras{kamambero sūriais} \textsuperscript{(10)}, \VUgras{griujero ratais}
\textsuperscript{(11)} ir šviežiais \VUgras{sviesto gabalais} \textsuperscript{(12)}.

%%K t15-03-1 p
3. --- Priešais ją \VUgras{daržovių pardavėja} \textsuperscript{(13)} siūlo savo
klientėms gražius \VUgras{pomidorus} \textsuperscript{(14)}, \VUgras{žiedinius
kopūstus} \textsuperscript{(15)}, \VUgras{salotas} \textsuperscript{(16)},
\VUgras{gūžinius kopūstus} \textsuperscript{(17)}, \VUgras{moliūgėlius}
\textsuperscript{(19)}, \VUgras{melionus} \textsuperscript{(18)} ir net
\VUgras{grybus} \textsuperscript{(20)}. Bet \VUgras{alaus išvežiotojas,} kuris
sustabdė savo \VUgras{pristatymo vežimą} \textsuperscript{(21)}, prikrautą
\VUgras{alaus statinių} \textsuperscript{(22)}, kaip tik prieš jos prekystalį, trukdo
jos klientams prieiti. Sodininkė atneša jai krepšį \VUgras{artišokų}
\textsuperscript{(23)}.

%%K t15-tit-2 sub
\VUsaut{4.0mm}
\VUpk{10.2pt}{40}{Pardavimas ir pirkimas.}
\VUsaut{3.0mm}

%%K t15-04-1 p
4. --- Turguje didelis judėjimas, nes atėjo \VUgras{aukciono} \textsuperscript{(24)}
valanda. \VUgras{Šeimininkės} \textsuperscript{(26)}, kurios ateina čia savo pirkinių,
eidamos sustoja prieš \VUgras{vidurių pardavėjos} \textsuperscript{(25)} būdelę,
kad nupirktų \VUgras{skrandžio vidučių} arba \VUgras{veršio galvą.}

%%K t15-05-1 p
5. --- Riebus \VUgras{virėjas} \textsuperscript{(27)} derasi dėl labai gražaus
\VUgras{tuno} pas \VUgras{žuvų pardavėją} \textsuperscript{(28)}, kuri savo nuožiūra
kelia savo kainas. Ji gavo didelį kiekį \VUgras{ungurių, jūrų liežuvių, upėtakių} ir
\VUgras{karpių.} Jos \VUgras{menkė} ir \VUgras{midijos} labai šviežios.

%%K t15-tit-3 sub
\VUsaut{4.0mm}
\VUpk{10.2pt}{40}{Pigios ir brangios prekės.}
\VUsaut{3.0mm}

%%K t15-06-1 p
6. --- \VUgras{Paukštienos pardavėja} \textsuperscript{(29)}, įsitaisiusi šalia jos,
labai sąžininga. Kai ji parduoda \VUgras{kalakutą, žąsį, antį} ar paprastą
\VUgras{viščiuką,} ji prašo tik teisingos kainos.

%%K t15-07-1 p
7. --- Aikštės kampe, antrajame aukšte, neseniai įsikūrė \VUgras{audinių
pardavėjas} \textsuperscript{(30)}, pas kurį pirkėjos visada tikros rasiančios labai
gražų \VUgras{staltiesių, servetėlių, prijuosčių} ir \VUgras{rankšluosčių} pasirinkimą.
Tame pačiame aukšte \VUgras{peilininkas} \textsuperscript{(31)} įsirengė savo dirbtuvę.
Jis galanda halės pardavėjų \VUgras{peilius} ir gamina sodininkams \VUgras{pjautuvus}
iš kiečiausio \VUgras{plieno.}

%%K t15-tit-4 sub
\VUsaut{4.0mm}
\VUpk{10.2pt}{40}{Prieskonių prekyba.}
\VUsaut{3.0mm}

%%K t15-08-1 p
8. --- Po jo dirbtuve neseniai atidaryta didelė maisto prekių krautuvė. Ji jau
nebėra ta paprasta ir nesvarbi senoji \VUgras{prieskonių krautuvė}
\textsuperscript{(32)}, kuri pardavinėjo tik įprastas prekes --- \VUgras{arbatą}
\textsuperscript{(33)}, \VUgras{šokoladą} \textsuperscript{(34)}, \VUgras{gabalinį
cukrų} \textsuperscript{(37)} ir \VUgras{muilą} \textsuperscript{(38)}. Ten parduodama
bet kokia prekė: \VUgras{traškūs sausainiai,} \VUgras{konservai}
\textsuperscript{(35)}, \VUgras{likeriai} \textsuperscript{(36)}, \VUgras{romas,}
\VUgras{konjakas, kiršas, anyžinis likeris} ir \VUgras{kiurasao.} Ten randama
žvėrienos: \VUgras{kiškis} \textsuperscript{(39)}, \VUgras{strazdai}
\textsuperscript{(40)}, \VUgras{kurapkos} \textsuperscript{(41)}, \VUgras{putpelės}
\textsuperscript{(42)}, \VUgras{laukinė antis} \textsuperscript{(43)} ar
\VUgras{fazanas} \textsuperscript{(44)}, taip pat lengvai kaip egzotiškų vaisių, tokių
kaip \VUgras{bananai} \textsuperscript{(46)} ir \VUgras{ananasai.}

%%K t15-09-1 p
9. --- Labai paslaugus prekybininko \VUgras{pameistrys} \textsuperscript{(45)}
pasirengęs duoti, ko norima: \VUgras{uogienės} \textsuperscript{(47)} --- serbentų ar
svarainių ---, \VUgras{taukų} \textsuperscript{(48)}, \VUgras{agurkėlių}
\textsuperscript{(49)} ar sūdytų \VUgras{sardinių} \textsuperscript{(50)}.

%%K t15-09-2 p
Tai labai patogu \VUgras{klientėms} \textsuperscript{(51)}, kurios randa, greta
įprasčiausių prekių, rečiausius pirmavaisius ir skaniausius vaisius: sultingas
\VUgras{kriaušes} \textsuperscript{(52)}, \VUgras{slyvas} \textsuperscript{(53)},
\VUgras{vynuoges} \textsuperscript{(54)}, \VUgras{persikus} \textsuperscript{(55)},
\VUgras{migdolus} \textsuperscript{(56)} ir \VUgras{riešutus} \textsuperscript{(57)}.

%%K t15-tit-5 sub
\VUsaut{4.0mm}
\VUpk{10.2pt}{40}{Klientai.}
\VUsaut{3.0mm}

%%K t15-10-1 p
10. --- \VUgras{Ryžiai} \textsuperscript{(58)} ir \VUgras{lęšiai}
\textsuperscript{(59)} yra prie rūkytų \VUgras{silkių} \textsuperscript{(60)} dėžės;
\VUgras{bulvės} \textsuperscript{(61)} ir \VUgras{pupelės} \textsuperscript{(63)} yra
šalia \VUgras{obuolių} \textsuperscript{(62)}, \VUgras{figų} \textsuperscript{(64)},
\VUgras{džiovintų slyvų} \textsuperscript{(65)} ir \VUgras{lazdyno riešutų}
\textsuperscript{(66)}. Toje krautuvėje randama šviežių \VUgras{kiaušinių}
\textsuperscript{(67)} ir \VUgras{alyvuogių} \textsuperscript{(68)}; todėl
\VUgras{krepšiai} \textsuperscript{(70)} ir \VUgras{pirkinių tinkleliai}
\textsuperscript{(73)} labai greitai prisipildo.

%%K t15-11-1 p
11. --- Tos puikios firmos \VUgras{kava} \textsuperscript{(72)} įgijo tarp
\VUgras{tarnaičių} \textsuperscript{(69)} ir \VUgras{virėjų} teisingą vardą, nes
senasis \VUgras{prekybininkas} \textsuperscript{(71)} pats rūpestingiausiai ją skrudina.

%%K t15-tit-6 sub
\VUsaut{4.0mm}
\VUpk{10.2pt}{40}{Reginiai.}
\VUsaut{3.0mm}

%%K t15-12-1 p
12. --- Ne visi eina į turgų apsirūpinti. Vaizdingame judėjime gatvelėje, kuri veda
į halę, matyti patys įvairiausi žmonės. Šalia malonaus \VUgras{mėsininko berno}
\textsuperscript{(75)}, kuris stumia prieš save \VUgras{karutį} \textsuperscript{(74)},
štai du atostogaujantys kareiviai, labai išdidūs galėdami parodyti savo blizgančią
uniformą: \VUgras{husaras} \textsuperscript{(76)} su mėlynu \VUgras{dolmanu} ir
\VUgras{plunksnuotu šaku,} ir \VUgras{zuavų kapralas} su pūstomis kelnėmis ir galva,
uždengta \VUgras{fesu.}

%%K t15-13-1 p
13. --- \VUgras{Kepėjas} \textsuperscript{(80)}, kuris stovi ant \VUgras{kepyklos}
\textsuperscript{(78)} slenksčio, ką tik išminkė savo \VUgras{duonos}
\textsuperscript{(79)} tešlą ir išėmė iš krosnies \VUgras{raguolius}
\textsuperscript{(81)}, \VUgras{bandeles} \textsuperscript{(82)} ir
\VUgras{paplotėlius} \textsuperscript{(83)}, kurie yra vitrinoje.

%%K t15-14-1 p
14. --- Virš kepyklos gyvena sumani \VUgras{baltinių siuvėja} \textsuperscript{(84)},
kuri samdo kelias \VUgras{siuvinėtojas} \textsuperscript{(85)}. Šalia jos gyvena
\VUgras{skalbėja-lygintoja} \textsuperscript{(86)}, kuri krakmolija
\VUgras{skalbinius} \textsuperscript{(88)}, juos išplovusi ir išdžiovinusi, ir kuri
lygina juos \VUgras{lygintuvu} \textsuperscript{(87)}.

%%K t15-tit-7 sub
\VUsaut{4.0mm}
\VUpk{10.2pt}{40}{Mėsa, žuvys ir daržovės.}
\VUsaut{3.0mm}

%%K t15-15-1 p
15. --- \VUgras{Mėsinėje} \textsuperscript{(89)}, esančioje pirmajame aukšte,
parduodama visokių \VUgras{mėsų: aviena} \textsuperscript{(90)}, \VUgras{jautiena}
\textsuperscript{(94)} ar \VUgras{veršiena} \textsuperscript{(92)} --- kaip
\VUgras{avienos kumpiai, bifšteksai, kepsniai} ir \VUgras{kotletai.}
\VUgras{Mėsininko pameistrys} \textsuperscript{(93)} išpjausto visas tas mėsas ir
atskirai deda \VUgras{smegeninius kaulus} \textsuperscript{(96)}. Prie mėsinės durų yra
\VUgras{austrių pardavėja} \textsuperscript{(97)}, kuri parduoda \VUgras{austres}
\textsuperscript{(98)}. Ji jas atidaro, jei to pageidaujama.

%%K t15-16-1 p
16. --- Visi tie prekybininkai moka patentą ar vietos mokestį; todėl
\VUgras{policininkas} \textsuperscript{(100)} negailestingai persekioja vargšes
\VUgras{daržovių pardavinėtojas} \textsuperscript{(99)}, kurios su jomis konkuruoja,
išvežiodamos savo vežimėliais \VUgras{morkas, ridikėlius, ropes, poringus} ar
\VUgras{smidrų glėbelius, šviežius žirnius, špinatus, rūgštynes, pipirnes} ir
\VUgras{petražoles.}

%%K t15-tit-8 sub
\VUsaut{4.0mm}
\VUpk{10.2pt}{40}{Saugokis policininko!}
\VUsaut{3.0mm}

%%K t15-17-1 p
17. --- Berniukas, kuris parduoda \VUgras{česnakus} \textsuperscript{(101)} ir
\VUgras{svogūnus,} labai gerai žino, kad ir jam negalima pardavinėti savo prekių prie
turgaus. Jis bėga, rizikuodamas apversti \VUgras{krabų, vėžių} ir \VUgras{krevečių}
krepšį, priklausantį \VUgras{langustų} ir \VUgras{omarų} \VUgras{pardavėjui}
\textsuperscript{(102)}.

%%K t15-18-1 p
18. --- Visi juda ir triukšmauja; bet tai netrukdo aiškiai girdėti klajojančio
\VUgras{stikliaus} \textsuperscript{(103)} balso nei \VUgras{anglininko}
\textsuperscript{(104)} šauksmo, kuris kelioms akimirkoms paliko savo vežimėlį, kad
pristatytų \VUgras{anglių maišelį} \textsuperscript{(105)}.

\VUsaut{6.0mm}
\VUfilet{22mm}
